\subsection{Floating-Point Data Formats}
\manualtablecaption{Floating-Point Scalar Sizes}
\begin{manualtabular}{@{}>{\raggedright\arraybackslash}p{0.45in}>{\raggedright\arraybackslash}p{0.55in}>{\raggedright\arraybackslash}p{0.45in}>{\raggedright\arraybackslash}p{0.50in}>{\raggedright\arraybackslash}p{1.55in}@{}}
\toprule
\rowcolor{ManualHeaderFill}
\textbf{Code} & \textbf{Suffix} & \textbf{Bits} & \textbf{Bytes} & \textbf{Name}\\
\midrule
\texttt{H} & \texttt{.H} & 16 & 2 & Half (conversion and storage)\\
\texttt{S} & \texttt{.S} & 32 & 4 & Single\\
\texttt{D} & \texttt{.D} & 64 & 8 & Double\\
\bottomrule
\end{manualtabular}

Floating-point scalar formats are little-endian in memory as well. H is IEEE-754 binary16 with 11 bits of
significand precision, minimum normal exponent $-14$, and minimum subnormal quantum $2^{-24}$. H is available to
the scalar instruction set only as a conversion and storage format; scalar arithmetic remains S/D. S is IEEE-754 binary32 with 24 bits of
significand precision, minimum normal exponent $-126$, and minimum subnormal quantum $2^{-149}$. D is IEEE-754
binary64 with 53 bits of significand precision, minimum normal exponent $-1022$, and minimum subnormal quantum
$2^{-1074}$.

\begin{BedrockListedFormatDiagram}{Floating-Point Register Encodings}[fig:floating-point-register-encodings]
\BedrockFormatRowRange{D / Fn[63:0]}{63}{0}{%
\BedrockFormatField{s}{1}
\BedrockFormatField{exponent}{11}
\BedrockFormatField{fraction}{52}
}
\BedrockFormatRowRange{S / Fn[63:0]}{63}{0}{%
\BedrockFormatField{zero}{32}
\BedrockFormatField{s}{1}
\BedrockFormatField{exponent}{8}
\BedrockFormatField{fraction}{23}
}
\end{BedrockListedFormatDiagram}
{\small
For a NaN, the most significant fraction bit is the quiet bit: bit 22 for S and bit 51 for D.
\texttt{s} denotes the sign bit. An S result occupies Fn bits 31..0 and writes zero to Fn bits 63..32.
\par}


The bit-level floating-point encoding, NaN policy, exception flags, and rounding behavior are defined by the
floating-point instruction and state descriptions; their memory byte order follows the same
least-significant-byte-first rule as integer data.
